\documentclass[11pt]{article}
\usepackage[T1]{fontenc}
\usepackage[margin=1in]{geometry}
\usepackage{fancyvrb}
\usepackage{xcolor}
\usepackage{enumitem}
\usepackage[hidelinks]{hyperref}

\definecolor{HardyBlue}{RGB}{24,64,112}
\setlength{\parindent}{0pt}
\setlength{\parskip}{0.7em}
\setlist[itemize]{leftmargin=1.5em}

\begin{document}

\begin{center}
{\LARGE\color{HardyBlue}\bfseries The Sum of the Square Roots of 2 and 3 Is Irrational\par}
\vspace{1em}
Formal status: Kernel verified\\
Faithfulness: User approved; independently reviewed by claude-opus-5 on claude (agreed)\\
Informal completeness: Not independently assessed\\
Document status: TeX compiled
\end{center}

\hrule
\vspace{1em}

\section*{Theorem}
For every rational number q, the real number q (viewed in ℝ via the canonical embedding ℚ ↪ ℝ) is not equal to √2 + √3. Equivalently, √2 + √3 is irrational.

In Lean: theorem irrational\_sqrt\_two\_add\_sqrt\_three : ∀ q : ℚ, (q : ℝ) ≠ Real.sqrt 2 + Real.sqrt 3

\section*{Approved interpretation}
\begin{itemize}
\item ℝ — the real numbers, Mathlib's `Real`
\item ℚ — the rational numbers, Mathlib's `Rat`, mapped into ℝ by the canonical coercion `Rat.cast : ℚ → ℝ` (written `(q : ℝ)`)
\item `Real.sqrt : ℝ → ℝ` — the nonnegative square-root function on ℝ (value 0 on negative inputs, irrelevant here since 2 > 0 and 3 > 0)
\item The literals 2 and 3 elaborate as real numbers, forced by the argument type of `Real.sqrt`
\item One explicit universal quantifier inside the proposition: `∀ q : ℚ`
\item The statement is closed (no free variables, no binders in the signature)
\item None; the claim is unconditional.
\item Standard Mathlib coercion `ℚ → ℝ` is available as a `RatCast ℝ` instance.
\item `Real.sqrt 2` and `Real.sqrt 3` denote the principal nonnegative real square roots, so the sum is the intended real number ≈ 3.14626.
\item Repair attempt 2: replaced the Mathlib predicate `Irrational` with its unfolded meaning `∀ q : ℚ, (q : ℝ) ≠ Real.sqrt 2 + Real.sqrt 3`. `Irrational` lives in `Mathlib.Data.Real.Irrational` and is not available under a restricted or fixed import set; the unfolded form needs only `Real`, `Rat`, the `ℚ → ℝ` cast, and `Real.sqrt`, so it elaborates under a much weaker import prefix. This is the definitional content of `Irrational x` (`x ∉ Set.range ((↑) : ℚ → ℝ)`), so nothing is weakened or strengthened.
\item Also dropped the dependence on `Set.range` and on set-membership notation, which the previous form pulled in implicitly through `Irrational`.
\item Used the ∀-with-≠ form rather than `¬ ∃ q : ℚ, ...` to avoid any negation/parenthesization ambiguity when the signature is assembled; the two are logically equivalent.
\item `binders` is left empty and the quantifier is carried inside `proposition`, so the statement remains closed; `proposition` contains only a Prop-valued term, never a `theorem` keyword or colon.
\item Removed the explicit `(2 : ℝ)` / `(3 : ℝ)` ascriptions: `Real.sqrt`'s argument type already forces the literals to ℝ, and fewer ascriptions means fewer places for elaboration to fail.
\item 'sqrt(2)' and 'sqrt(3)' are read as principal (nonnegative) square roots, the standard convention; no other branch is asserted.
\item The claim is neither weakened (e.g. to 'not an integer') nor strengthened (e.g. to a degree-4 or algebraic-independence claim).
\end{itemize}

\section*{Human-readable proof}
Overview. The direct approach — arguing about √2 and √3 separately — is awkward, because a sum of two irrationals can perfectly well be rational (for example √2 + (1 − √2) = 1). The efficient route is to square the supposed rational value: squaring √2 + √3 collapses the two surds into the single surd √6, and the resulting identity exhibits √6 as a rational number. Since √6 is irrational, no such rational value can exist.

Setup. Suppose, for contradiction, that there is a rational number q with

    (q : ℝ) = √2 + √3.                                            (∗)

Here √ denotes the real square-root function, which on nonnegative arguments satisfies (√x)² = x.

Step 1: elementary square-root facts. Because 2 ≥ 0 and 3 ≥ 0 we have

    (√2)² = 2,    (√3)² = 3.

Also, because 2 ≥ 0, the multiplicativity of the square root on nonnegative arguments gives

    √2 · √3 = √(2 · 3) = √6.

(In the formal proof these three facts are supplied by Real.sq\_sqrt and Real.sqrt\_mul together with a numeric computation.)

Step 2: square the hypothesis. Squaring both sides of (∗) and expanding the binomial:

    q² = (√2 + √3)² = (√2)² + 2·√2·√3 + (√3)².

Substituting the three facts of Step 1 yields

    q² = 2 + 2√6 + 3 = 5 + 2√6.                                   (†)

Note the point of the manoeuvre: the two individual surds have disappeared, and exactly one surd, √6, remains, appearing linearly.

Step 3: √6 is forced to be rational. Solving (†) for √6 is pure linear algebra over ℝ:

    √6 = (q² − 5)/2.

Now let r := (q² − 5)/2, computed inside ℚ. This is a legitimate rational number, since ℚ is a field and q ∈ ℚ. Because the embedding ℚ ↪ ℝ is a ring homomorphism, it commutes with squaring, subtraction and division by 2, so the image of r in ℝ equals (q² − 5)/2 computed in ℝ, which is √6. Hence √6 lies in the range of the embedding ℚ ↪ ℝ — that is, √6 is rational. (Recall that in this formalization "Irrational x" is by definition the statement that x is not in the range of that embedding, so producing the witness r is exactly what is needed to contradict irrationality.)

Step 4: but √6 is irrational. The number 6 is a natural number that is not a perfect square: 2² = 4 < 6 < 9 = 3², so 6 is strictly between two consecutive squares and therefore is not the square of any natural number. For a natural number n that is not a perfect square, √n is irrational. (The classical argument: if √n = a/b in lowest terms, then a² = n b², and comparing the exponent of any prime dividing n to an odd power on the two sides gives a contradiction; equivalently, a rational root of the monic integer polynomial X² − n must be an integer by the rational root theorem, and no integer squares to 6.) In the formal development this step is discharged by norm\_num's irrational-square-root extension, which is exactly the "strictly between consecutive squares" criterion mechanized.

Step 5: contradiction. Steps 3 and 4 assert respectively that √6 is and is not in the range of ℚ ↪ ℝ. This is a contradiction, so the assumption (∗) is untenable.

Conclusion. No rational number q satisfies (q : ℝ) = √2 + √3; that is, √2 + √3 is irrational. □

Remarks.

1. Where the hypotheses were used. Nonnegativity of 2 and 3 was needed twice: to get (√x)² = x, and to get multiplicativity √2·√3 = √6. Both would fail without it, so the numeric side conditions are not decorative.

2. No case split on q = 0 is required. An alternative proof isolates √2 by writing √3 = q − √2 and squaring to obtain √2 = (q² − 1)/(2q), which needs the separate observation that q ≠ 0 (true since √2 + √3 > 0). The √6 route avoids this division entirely, which is why it is shorter both on paper and in Lean.

3. Generality. The same argument shows √m + √n is irrational whenever m, n ≥ 0 are integers with mn not a perfect square: squaring gives √(mn) = (q² − m − n)/2, and irrationality of √(mn) closes the case. The instance m = 2, n = 3 has mn = 6, which is not a perfect square.

Status of this exposition. This write-up is a human-readable rendering of the argument; it has not been independently assessed here. The corresponding Lean proof body was reported as accepted by the Lean checker, but a successful tool run is feedback only — acceptance of the solution rests solely with Hardy's FinalVerifier.

\section*{Known gaps}
\begin{itemize}
\item This informal exposition has not been independently assessed. It is a human-readable rendering of the formal argument, not a verified object in its own right; any discrepancy between this prose and the Lean proof body is possible in principle.
\item A successful Lean tool run is feedback only. Only Hardy's FinalVerifier can accept the solution, and no such acceptance is claimed here.
\item Step 4 (irrationality of √6) is not proved from first principles in the formal development: it is discharged by Mathlib's norm\_num irrational-square-root extension, which mechanizes the criterion 'k² < n < (k+1)² implies √n is irrational'. The classical justification sketched in the prose (rational root theorem / prime-exponent parity) is offered as explanation, not as the formally checked route, and the correctness of that extension is taken on trust as part of the pinned Mathlib.
\item Standard facts used without separate proof, as library lemmas of the pinned Mathlib: (√x)² = x for x ≥ 0 (Real.sq\_sqrt), √(xy) = √x·√y for x ≥ 0 (Real.sqrt\_mul), the ring-homomorphism property of the coercion ℚ → ℝ, and the definitional unfolding of Irrational as non-membership in the range of that coercion.
\item Steps 2 and 3 are carried out in the formal proof by the automation tactics nlinarith, push\_cast and linarith rather than by explicitly written algebraic manipulations; the human-readable expansion of the binomial and the linear solve given above reconstruct what those tactics are expected to establish, but the exact certificates they produced were not inspected.
\item The generalization stated in Remark 3 (√m + √n irrational when mn is not a perfect square) is asserted as context only. It was not formalized or checked, and lies outside the Frozen Claim.
\end{itemize}

\section*{Exact Lean statement}
\begin{Verbatim}[fontsize=\small]
theorem irrational_sqrt_two_add_sqrt_three : ∀ q : ℚ, (q : ℝ) ≠ Real.sqrt 2 + Real.sqrt 3
\end{Verbatim}

\section*{Verification appendix}
\textbf{Axioms used:} propext, Classical.choice, Quot.sound

\begin{Verbatim}[fontsize=\small]
Frozen Claim SHA-256: 67bbf6be3f6cce3fbcf0c18652aeedf6f61ed5d0c09a91ebedbfd36aa4077204
Lean source SHA-256: 7280e5b6de3e27f93d0d41605a25aafda4c36fb3d8955945633cd93c5cbebec3
Verification SHA-256: 35f53ddbf64ff223571daedcd425260158466ce1e30ff2a44144b13088e4b2f9
Prompt set SHA-256: 50c207a708e3696c9846957a60089ccd0fb4a6ddbc6a6920449887455cf3f5e2
Tectonic bundle SHA-256: 6ffe055852f8faf66c0acbe1a7fb27f87b869a90bad1204f3bf4d9683f597c7c
\end{Verbatim}

\begin{samepage}
\textbf{Run and tool identities}
\begin{itemize}
\item Run ID: 3c689084-8c38-41f6-ad81-c3bf4130b0ce
\item Model: claude-opus-5
\item Backend: claude
\item Runtime SDK: 0.2.127
\item Lean: 4.33.1
\item Mathlib: 0df444a360eaa60ab8c11dca51a86af692955474
\item Tectonic: 0.16.9
\item Tectonic bundle: https://data1.fullyjustified.net/tlextras-2022.0r0.tar
\end{itemize}
\end{samepage}

\vfill
\small
Lean kernel verification, faithfulness (approved by the user and read back by an
independent model before proof search), heuristic informal review, and document
compilation are independent statuses. This trusted-development MVP
does not sandbox generated Lean or TeX.

\end{document}
